\documentclass[12pt,a4paper]{article}
\usepackage[margin=25mm]{geometry}
\usepackage{amsmath,amssymb,amsthm}
\usepackage[T1]{fontenc}
\usepackage{lmodern}
\setlength{\parindent}{0pt}
\setlength{\parskip}{5pt plus 1pt}
\newtheorem{lemma}{Lemma}
\newcommand{\pat}[1]{\boldsymbol{#1}}
\begin{document}
\begin{center}
{\Large\bfseries A paper-style proof of the\\[3pt]
Lean formalization of Lemma 4}
\end{center}
\smallskip

\setcounter{lemma}{3}
\begin{lemma}
For any sequence of PPS, VPS, PPM operations, there exists an equivalent
cascade that consists of a sequence of PPS operations, followed by a
sequence of VPS operations, and finally a sequence of PPM operations.
\end{lemma}

\begin{proof}
Let $M$ be a singular conditional probability matrix with finite alphabets
and at least one input symbol, and let $N$ be obtained from $M$ by a finite
cascade of PPS, VPS and PPM operations. We work with active columns, meaning
that no column is identically zero, and regard output-column permutations
as relabelings. Write the columns of $M$ as $c_i\pat{p}_i$, where $c_i>0$
and $\pat{p}_i$ is a nonzero binary vector. We construct an equivalent
cascade in the required order by keeping track of how pieces of these
original columns contribute to each intermediate matrix.

Each allowed operation preserves singularity and activity. Hence the
columns of every intermediate matrix $Q$ can be written as
$d_j\pat{q}_j$, with $d_j>0$ and nonzero binary patterns $\pat{q}_j$.
For every original column $i$, we will maintain a finite family of
positive weights $w_{ia}$ summing to one. For each weighted copy $a$, we
will maintain a finite family of nonzero, pairwise disjoint binary
patterns $\pat{p}_{iab}$ whose sum is $\pat{p}_i$. The index range of $a$
may depend on $i$, and that of $b$ on $(i,a)$. Thus
\begin{equation}
  \sum_a w_{ia}=1, \qquad \sum_b \pat{p}_{iab}=\pat{p}_i.
  \label{eq:partitions}
\end{equation}
Assign each piece $c_iw_{ia}\pat{p}_{iab}$ a destination column
$\ell(i,a,b)$ of $Q$, with the properties
\begin{equation}
  \pat{p}_{iab}=\pat{q}_{\ell(i,a,b)}, \qquad
  Q[:,j]=\sum_{\ell(i,a,b)=j} c_iw_{ia}\pat{p}_{iab}.
  \label{eq:destinations}
\end{equation}
In other words, every piece has exactly the pattern of its destination,
and the pieces assigned to a column sum to that column. We also retain
finite binary trees recording the proportional splits that produce the
weights and the support splits that produce the pattern partitions.

We establish this description by induction on the length of the original
cascade. For $Q=M$, take one copy of weight one and one pattern
$\pat{p}_i$ for each column $i$, with destination $i$. Suppose now that
the description holds for $Q$, and consider the next operation.

\emph{PPS.} Suppose column $l$ is split in the ratio
$\alpha:(1-\alpha)$, where $0<\alpha<1$. Replace \emph{every} weighted
copy of every original column by two copies with weights $\alpha w_{ia}$
and $(1-\alpha)w_{ia}$, retaining the same pattern partition in each.
For pieces with destination $l$, send the two copies to the two new
columns, respectively. For every other piece, send both copies to its
unchanged destination. Its two contributions add to the original
contribution, whereas the contributions to column $l$ are split in the
required ratio. The support identities are unchanged, and all weights
remain positive and sum to one. The new weight tree consists of a split
with ratio $\alpha$ followed by a copy of the old tree on each branch.

\newpage
\emph{VPS.} Suppose column $l$ is split into two nonzero, disjoint
patterns $\pat{q}_{l,1}$ and $\pat{q}_{l,2}$. Every piece assigned to $l$
has pattern $\pat{q}_l$. Replace its pattern by these same two patterns
and assign the resulting pieces to the corresponding new columns,
leaving its weight unchanged. Both pieces are nonzero because the given
VPS is legal. All other pieces are unchanged. This refines the relevant
pattern partitions, preserves the column-sum identities, and adds one
binary support split at each affected leaf.

\emph{PPM.} Suppose two columns with the same pattern are merged. Leave
all weights and pattern partitions unchanged, and redirect the pieces
assigned to either column to the merged column. Their patterns still
agree with that of their destination, and their total contribution is
the sum of the two original columns. Hence the required description is
preserved.

Applying the induction to the whole cascade gives such a description of
$N$. We now realize it in three stages. First, perform the recorded
proportional splits on each original column, obtaining the columns
$c_iw_{ia}\pat{p}_i$. Second, perform the recorded support splits on each
weighted copy, obtaining the pieces $c_iw_{ia}\pat{p}_{iab}$. Every
support split is legal: both branches contain nonzero parts of the
corresponding partition. The first stage therefore uses only PPS
operations, and the second only VPS operations.

Finally, for each column $j$ of $N$, merge all pieces whose destination
is $j$. These pieces are positive multiples of the same nonzero pattern,
so they can be merged successively using PPM operations. Every target
column receives at least one piece, since it is nonzero and equals the
sum of its assigned pieces. The column-sum identity
in~\eqref{eq:destinations} shows that the resulting column is precisely
$N[:,j]$. Thus the endpoint is $N$ up to a permutation of its output
columns.

The construction is finite: the description is built by induction on
the given finite cascade, and each of its recorded binary trees is
finite. Traversing these trees and merging the finite destination groups
produces a finite PPS--VPS--PPM cascade; any of the three blocks may be
empty. This proves the lemma.
\end{proof}
\end{document}
